% ============================================================================
\section{GPAS Micro-Batch Allocation}
\label{sec:theory}
% ============================================================================

\subsection{Noise in optimizer coordinates}

AdamW's current second-moment state rescales coordinate $k$ by
$D_k=(\sqrt{v_k}+\epsilon)^{-1}$. We therefore measure the
one-micro-batch noise of task $i$ in the same coordinates. Assuming independent
micro-batches,
\begin{equation}
\begin{split}
    e_i&=\mathbb E\bigl\|D(g_i-\mathbb E[g_i])\bigr\|_2^2,\\
    V(m)&=\mathbb E\bigl\|D(A-\nabla F)\bigr\|_2^2
          =\sum_{i=1}^M\frac{w_i^2e_i}{m_i}.
\end{split}
\label{eq:batch_variance}
\end{equation}
Because every task is present, $V(m)$ measures within-task sampling noise and
does not include task-selection noise. For a fixed student state, step size,
and $D$, the allocation affects the one-step descent bound only through $V(m)$
(Appendix~\ref{app:descent}).

Minimizing $V(m)$ for a fixed budget of $G$ micro-batches gives
\begin{equation}
    m_i\ \propto\ w_i\sqrt{e_i}.
    \label{eq:batch_gpas_allocation}
\end{equation}
We call this rule Gradient-Preconditioned Adaptive Sampling (GPAS). It is
Neyman allocation---the classical rule that samples each stratum in proportion
to its standard deviation---in the coordinates that AdamW uses. We clip the
continuous counts to $[m_{\min},m_{\max}]$, rescale the remaining counts to sum
to $G$, and round by largest remainder. Appendix~\ref{app:allocation_proofs}
gives the short derivation.

\subsection{Cost-GPAS}

Long-response tasks can cost more per micro-batch, so minimizing variance per
step need not minimize variance per hour. Let $\tau_i$ be the measured marginal
time of one task-$i$ micro-batch. Let $C$ denote fixed per-step time for
synchronization, teacher transfer, and the optimizer update. Under the serial
teacher layout used here, Cost-GPAS minimizes
\begin{equation}
    J_C(m)=\left(C+\sum_{i=1}^M m_i\tau_i\right)
           \left(\sum_{i=1}^M\frac{w_i^2e_i}{m_i}\right).
    \label{eq:batch_cost_objective}
\end{equation}
The first factor is the measured step-time model and the second is $V(m)$; the
product is proportional to the variance of a mean gradient estimate under a
fixed wall-clock budget. With no fixed overhead, the solution is
\[
    m_i\ \propto\ w_i\sqrt{e_i/\tau_i}.
\]
When $C$ is large, the solution approaches GPAS. For the measured $C>0$ case,
we enumerate the small feasible set of four-task integer allocations. The first
step uses Uniform because it has no preceding noise or timing observation.

\subsection{Online estimates and variance ratio}

All statistics for the next allocation come from the completed step. For each
task, $\widehat e_i$ is the unbiased sample variance of the vectors $Dg_{i,s}$,
summed over coordinates; equivalently, its numerator is
$\sum_s\|Dg_{i,s}\|_2^2-m_i\|D\bar g_i\|_2^2$ and its denominator is
$m_i-1$. The minimum count of two makes this estimate available on every step.
The default uses the current estimate without smoothing.
\missing{Report the observed step-to-step change of $\widehat e_i$ and revise
this choice if the four-teacher trace is unstable.}

Response lengths make timings more variable, so $\tau_i$ and $C$ use
exponential averages with decay $0.9$. We measure the predicted variance change
relative to the bounded Uniform allocation as
\begin{equation}
    H=\frac{V(m^{\rm unif})}{V(m)}.
    \label{eq:allocation_variance_ratio}
\end{equation}
When $H$ is close to one, the predicted variance under adaptive counts is close
to the Uniform value. With $m_i^{\rm unif}=4$ and $m_i\leq8$, $H\leq2$. For
Cost-GPAS, $H$ measures variance only; end-to-end GPU hours measure time
efficiency.

\subsection{Implementation and optional second-moment correction}

GPAS works with conventional AdamW. Its second-moment observation $A^2$
contains the allocation-dependent noise term
$\sum_i w_i^2\operatorname{Var}(g_i)/m_i$. As an optional consistency
refinement, we form
\begin{equation}
    U=\sum_{i=1}^M\frac{w_i}{m_i}\sum_{s=1}^{m_i}g_{i,s}^2,
    \qquad v_{\rm obs}=U/G.
    \label{eq:taskwise_square_observation}
\end{equation}
The expectation of $U/G$ is independent of $m$. Its noise scale matches the
conventional observation at proportional allocation $m_i=Gw_i$, avoiding the
roughly $\sqrt G=4\times$ effective-learning-rate change that would result from
using $U$ itself in noise-dominated coordinates. The remaining difference in
the signal term is derived in Appendix~\ref{app:moment_target_proof}. We call
this optional variant taskwise AdamW (TW) and test it separately from the main
allocation comparison.

\begin{algorithm}[H]
\caption{One optimizer step and the next GPAS allocation.}
\label{alg:batch_gpas}
\small
\begin{minipage}{0.97\linewidth}
\begin{enumerate}
    \item Use the previous step's $e_i$ (and $\tau_i,C$ for Cost-GPAS) to
    choose bounded integer counts $m$; use Uniform on the first step.
    \item For every task $i$, take $m_i b$ new prompts, sample one response per
    prompt, obtain token-level teacher feedback, and accumulate each
    micro-batch gradient with weight $w_i/m_i$.
    \item From the same gradients and execution trace, estimate $e_i$, update
    the timing averages, and log counts, losses, tokens, $H$, and timing.
    \item Apply one conventional AdamW update, or the optional TW update using
    $A$ for the first moment and $U/G$ for the second moment.
\end{enumerate}
\end{minipage}
\end{algorithm}
